\section{Confidential I/O protocol、trusted device interface、network endpoint}

\begin{frame}[t,shrink=6]{Confidential I/O protocol、trusted device interface、network endpoint：方向开场}
\DeckKicker{方向开场}
\SlideMeta{证据等级}{方向综述}{来源状态：公开材料综合}
{\large\bfseries\color{Ink} Confidential I/O protocol、trusted device interface、network endpoint 的核心是先界定保护对象、管理边界和证据强度，再用三篇主讲材料串起技术演进。\par}\vspace{1.2mm}
\begin{columns}[T,totalwidth=\textwidth]
  \begin{column}{0.45\textwidth}
    \fcolorbox{Rule}{Paper}{%
  \begin{minipage}[t]{0.97\linewidth}
    \vspace{1.1mm}
    \hspace{1.4mm}\begin{minipage}{0.92\linewidth}
      \PanelTitle{内容说明}
      \scriptsize \begin{itemize}
  \item 把 SPDM/TDISP/PCIe IDE/attested TLS 与 trusted device lifecycle 组合起来写。
  \item 固定三篇主讲材料；TLS+RA 作为 endpoint/channel-binding 主证据，FOLIO、SEV-TIO、TDISP/TDISP XT 和 BlueField DPU evidence 为 auxiliary/contrast。
  \item 先讲方向痛点和设计轴，再逐篇说明贡献、限制、证据强度和可落地条件。
  \item 对规范、Survey、draft、vendor evidence 保持显式标注，避免把背景材料写成一手系统证明。
\end{itemize}
    \end{minipage}
    \vspace{1mm}
  \end{minipage}}
  \end{column}
  \begin{column}{0.49\textwidth}
    \fcolorbox{Rule}{Panel}{%
  \begin{minipage}[t]{0.97\linewidth}
    \vspace{1.1mm}
    \hspace{1.4mm}\begin{minipage}{0.92\linewidth}
      \PanelTitle{三篇主讲材料的分工}
      \scriptsize {\tiny
\begin{tabularx}{\linewidth}{@{}YYY@{}}
\toprule
\textbf{论文} & \textbf{定位} & \textbf{证据边界} \\
\midrule
SPDM device identity and measurement & 开创/基础入口 & E0 / Spec-standard SOTA \\
CoVE-IO and TDISP trusted-device lifecycle & 代表性改进 & E3 / 草案/未批准规范 \\
Separate but Together: TLS channel binding with remote... & 当前 SOTA/补充边界 & E1 / Peer-reviewed endpoint channel-binding evidence \\
\bottomrule
\end{tabularx}
}
    \end{minipage}
    \vspace{1mm}
  \end{minipage}}
  \end{column}
\end{columns}
\SourceFootnote{来源：论文原文、官方规范或公开材料｜证据等级：见本页 badge}
\end{frame}


\begin{frame}[t,shrink=6]{SPDM device identity and measurement：内容摘要}
\DeckKicker{内容摘要}
\SlideMeta{证据等级}{E0 / Spec-standard SOTA}{来源状态：本地 PDF 已核验}
{\large\bfseries\color{Ink} SPDM 定义 device/component identity、measurement、capability negotiation、certificate/challenge 和 key exchange\par}\vspace{1.2mm}
\begin{columns}[T,totalwidth=\textwidth]
  \begin{column}{0.45\textwidth}
    \fcolorbox{Rule}{Paper}{%
  \begin{minipage}[t]{0.97\linewidth}
    \vspace{1.1mm}
    \hspace{1.4mm}\begin{minipage}{0.92\linewidth}
      \PanelTitle{内容说明}
      \scriptsize \begin{itemize}
  \item SPDM 定义 device/component identity、measurement、capability negotiation...
  \item SPDM is the first-party specification foundation for device identity, measurement, and...
\end{itemize}
    \end{minipage}
    \vspace{1mm}
  \end{minipage}}
  \end{column}
  \begin{column}{0.49\textwidth}
    \fcolorbox{Rule}{Panel}{%
  \begin{minipage}[t]{0.97\linewidth}
    \vspace{1.1mm}
    \hspace{1.4mm}\begin{minipage}{0.92\linewidth}
      \PanelTitle{证据定位}
      \scriptsize {\tiny
\begin{tabularx}{\linewidth}{@{}YY@{}}
\toprule
\textbf{字段} & \textbf{内容} \\
\midrule
论文定位 & 开创/基础入口 \\
材料类型 & 规范/标准材料 \\
证据等级 & E0 / Spec-standard SOTA \\
来源状态 & 本地 PDF 已核验 \\
\bottomrule
\end{tabularx}
}
    \end{minipage}
    \vspace{1mm}
  \end{minipage}}
  \end{column}
\end{columns}
\SourceFootnote{来源：SPDM device identity and measurement｜证据等级：E0 / Spec-standard SOTA}
\end{frame}


\begin{frame}[t,shrink=6]{SPDM device identity and measurement：研究背景}
\DeckKicker{研究背景}
\SlideMeta{证据等级}{E0 / Spec-standard SOTA}{来源状态：本地 PDF 已核验}
{\large\bfseries\color{Ink} Confidential workload 使用真实设备时，不能只相信 host 枚举出的 PCIe requester ID、VF 配置或驱动状态\par}\vspace{1.2mm}
\begin{columns}[T,totalwidth=\textwidth]
  \begin{column}{0.45\textwidth}
    \fcolorbox{Rule}{Paper}{%
  \begin{minipage}[t]{0.97\linewidth}
    \vspace{1.1mm}
    \hspace{1.4mm}\begin{minipage}{0.92\linewidth}
      \PanelTitle{内容说明}
      \scriptsize \begin{itemize}
  \item Confidential workload 使用真实设备时，不能只相信 host 枚举出的 PCIe requester ID、VF 配置或驱动状态
\end{itemize}
    \end{minipage}
    \vspace{1mm}
  \end{minipage}}
  \end{column}
  \begin{column}{0.49\textwidth}
    \fcolorbox{Rule}{Panel}{%
  \begin{minipage}[t]{0.97\linewidth}
    \vspace{1.1mm}
    \hspace{1.4mm}\begin{minipage}{0.92\linewidth}
      \PanelTitle{问题边界}
      \scriptsize {\tiny
\begin{tabularx}{\linewidth}{@{}YY@{}}
\toprule
\textbf{维度} & \textbf{说明} \\
\midrule
保护对象 & Confidential I/O protocol、trusted device interface、network endpoint \\
传统瓶颈 & 把 SPDM/TDISP/PCIe IDE/attested TLS 与 trusted device lifecycle 组合起来写。 \\
本文入口 & SPDM device identity and measurement \\
证据限制 & E0 / Spec-standard SOTA \\
\bottomrule
\end{tabularx}
}
    \end{minipage}
    \vspace{1mm}
  \end{minipage}}
  \end{column}
\end{columns}
\SourceFootnote{来源：SPDM device identity and measurement｜证据等级：E0 / Spec-standard SOTA}
\end{frame}


\begin{frame}[t,shrink=6]{SPDM device identity and measurement：解决方案}
\DeckKicker{解决方案}
\SlideMeta{证据等级}{E0 / Spec-standard SOTA}{来源状态：本地 PDF 已核验}
{\large\bfseries\color{Ink} 通过 requester/responder 协议、certificate chain、measurement blocks、transcript 和 secured session 建立 device trust...\par}\vspace{1.2mm}
\begin{columns}[T,totalwidth=\textwidth]
  \begin{column}{0.45\textwidth}
    \fcolorbox{Rule}{Paper}{%
  \begin{minipage}[t]{0.97\linewidth}
    \vspace{1.1mm}
    \hspace{1.4mm}\begin{minipage}{0.92\linewidth}
      \PanelTitle{内容说明}
      \scriptsize \begin{itemize}
  \item 通过 requester/responder 协议、certificate chain、measurement blocks、transcript 和 secured...
\end{itemize}
    \end{minipage}
    \vspace{1mm}
  \end{minipage}}
  \end{column}
  \begin{column}{0.49\textwidth}
    \fcolorbox{Rule}{Panel}{%
  \begin{minipage}[t]{0.97\linewidth}
    \vspace{1.1mm}
    \hspace{1.4mm}\begin{minipage}{0.92\linewidth}
      \PanelTitle{核心设计拆解}
      \scriptsize \begin{itemize}
  \item 01. Device certificate and measurement retrieval
  \item 02. Challenge-response and transcript binding
  \item 03. Key exchange for secured control messages
\end{itemize}
    \end{minipage}
    \vspace{1mm}
  \end{minipage}}
  \end{column}
\end{columns}
\SourceFootnote{来源：SPDM device identity and measurement｜证据等级：E0 / Spec-standard SOTA}
\end{frame}


\begin{frame}[t,shrink=6]{SPDM device identity and measurement：实验结果}
\DeckKicker{实验结果}
\SlideMeta{证据等级}{E0 / Spec-standard SOTA}{来源状态：本地 PDF 已核验}
{\large\bfseries\color{Ink} 规范/Survey，无新实验；SPDM device identity and measurement 只支撑术语、taxonomy、接口语义或证据边界。\par}\vspace{1.2mm}
\begin{columns}[T,totalwidth=\textwidth]
  \begin{column}{0.45\textwidth}
    \fcolorbox{Rule}{Paper}{%
  \begin{minipage}[t]{0.97\linewidth}
    \vspace{1.1mm}
    \hspace{1.4mm}\begin{minipage}{0.92\linewidth}
      \PanelTitle{内容说明}
      \scriptsize \begin{itemize}
  \item 本页不展示独立系统实验，也不把二手归纳写成一手结果。
  \item 可支撑：SPDM is the first-party specification foundation for device identity, measurement...
  \item 证据等级：E0 / Spec-standard SOTA；来源状态：本地 PDF 已核验。
  \item 涉及具体平台、协议状态或数值时，转到原始论文、官方规范或厂商材料核验。
\end{itemize}
    \end{minipage}
    \vspace{1mm}
  \end{minipage}}
  \end{column}
  \begin{column}{0.49\textwidth}
    \fcolorbox{Rule}{Panel}{%
  \begin{minipage}[t]{0.97\linewidth}
    \vspace{1.1mm}
    \hspace{1.4mm}\begin{minipage}{0.92\linewidth}
      \PanelTitle{实验/证据边界}
      \scriptsize {\tiny
\begin{tabularx}{\linewidth}{@{}YYY@{}}
\toprule
\textbf{对象} & \textbf{可支撑} & \textbf{边界} \\
\midrule
数据/来源 & SPDM device identity and measurement & 本地 PDF 已核验 \\
Baseline/对照 & 以论文或规范原文为准 & 不跨材料外推 \\
指标/对象 & 只使用当前来源可证明的信息 & E0 / Spec-standard SOTA \\
使用边界 & 可进入本方向演进图 & 不能替代更强证据 \\
\bottomrule
\end{tabularx}
}
    \end{minipage}
    \vspace{1mm}
  \end{minipage}}
  \end{column}
\end{columns}
\SourceFootnote{来源：SPDM device identity and measurement｜证据等级：E0 / Spec-standard SOTA}
\end{frame}


\begin{frame}[t,shrink=6]{SPDM device identity and measurement：文章评价}
\DeckKicker{文章评价}
\SlideMeta{证据等级}{E0 / Spec-standard SOTA}{来源状态：本地 PDF 已核验}
{\large\bfseries\color{Ink} 评价：SPDM device identity and measurement 适合做 taxonomy/规范入口；规范/Survey，无新实验，不能替代一手系统证据。\par}\vspace{1.2mm}
\begin{columns}[T,totalwidth=\textwidth]
  \begin{column}{0.45\textwidth}
    \fcolorbox{Rule}{Paper}{%
  \begin{minipage}[t]{0.97\linewidth}
    \vspace{1.1mm}
    \hspace{1.4mm}\begin{minipage}{0.92\linewidth}
      \PanelTitle{内容说明}
      \scriptsize \begin{itemize}
  \item 设计优势：是 CoVE-IO、TDISP、PCIe IDE keying 和 vendor device attestation 的共同 device-evidence 底座。
  \item 局限性：单独成功的 SPDM session 不能证明 TDI state、IOMMU/IOPMP policy、IDE link 或 device firmware 隔离正确。
  \item 商业化潜力：适合 fleet-scale device attestation；风险在证书供应链、measurement schema 和 firmware update...
  \item 规范/Survey，无新实验；具体平台结论转到一手材料核验。
\end{itemize}
    \end{minipage}
    \vspace{1mm}
  \end{minipage}}
  \end{column}
  \begin{column}{0.49\textwidth}
    \fcolorbox{Rule}{Panel}{%
  \begin{minipage}[t]{0.97\linewidth}
    \vspace{1.1mm}
    \hspace{1.4mm}\begin{minipage}{0.92\linewidth}
      \PanelTitle{文章评价}
      \scriptsize {\tiny
\begin{tabularx}{\linewidth}{@{}YY@{}}
\toprule
\textbf{维度} & \textbf{判断} \\
\midrule
设计优势 & 是 CoVE-IO、TDISP、PCIe IDE keying 和 vendor device attestation 的共同。 \\
主要局限 & 单独成功的 SPDM session 不能证明 TDI state、IOMMU/IOPMP policy、IDE link 或 device... \\
商业落地 & 适合 fleet-scale device attestation；风险在证书供应链、measurement schema 和 firmware... \\
讲稿定位 & 开创/基础入口; 证据强度 E0 \\
\bottomrule
\end{tabularx}
}
    \end{minipage}
    \vspace{1mm}
  \end{minipage}}
  \end{column}
\end{columns}
\SourceFootnote{来源：SPDM device identity and measurement｜证据等级：E0 / Spec-standard SOTA}
\end{frame}


\begin{frame}[t,shrink=6]{CoVE-IO and TDISP trusted-device lifecycle：内容摘要}
\DeckKicker{内容摘要}
\SlideMeta{证据等级}{E3 / 草案/未批准规范}{来源状态：本地 PDF 已核验}
{\large\bfseries\color{Ink} CoVE-IO 把 TVM trusted I/O 拆成 TDI/TDM/DSM、SPDM、TDISP、PCIe IDE、IOMMU、trusted MSI 和 CoVE ABI\par}\vspace{1.2mm}
\begin{columns}[T,totalwidth=\textwidth]
  \begin{column}{0.45\textwidth}
    \fcolorbox{Rule}{Paper}{%
  \begin{minipage}[t]{0.97\linewidth}
    \vspace{1.1mm}
    \hspace{1.4mm}\begin{minipage}{0.92\linewidth}
      \PanelTitle{内容说明}
      \scriptsize \begin{itemize}
  \item CoVE-IO 把 TVM trusted I/O 拆成 TDI/TDM/DSM、SPDM、TDISP、PCIe IDE、IOMMU、trusted MSI 和 CoVE ABI
  \item RISC-V CoVE-IO draft maps SPDM and TDISP-style ideas into the RISC-V TEE-I/O track.
\end{itemize}
    \end{minipage}
    \vspace{1mm}
  \end{minipage}}
  \end{column}
  \begin{column}{0.49\textwidth}
    \fcolorbox{Rule}{Panel}{%
  \begin{minipage}[t]{0.97\linewidth}
    \vspace{1.1mm}
    \hspace{1.4mm}\begin{minipage}{0.92\linewidth}
      \PanelTitle{证据定位}
      \scriptsize {\tiny
\begin{tabularx}{\linewidth}{@{}YY@{}}
\toprule
\textbf{字段} & \textbf{内容} \\
\midrule
论文定位 & 代表性改进 \\
材料类型 & 规范/标准材料 \\
证据等级 & E3 / 草案/未批准规范 \\
来源状态 & 本地 PDF 已核验 \\
\bottomrule
\end{tabularx}
}
    \end{minipage}
    \vspace{1mm}
  \end{minipage}}
  \end{column}
\end{columns}
\SourceFootnote{来源：CoVE-IO and TDISP trusted-device lifecycle｜证据等级：E3 / 草案/未批准规范}
\end{frame}


\begin{frame}[t,shrink=6]{CoVE-IO and TDISP trusted-device lifecycle：研究背景}
\DeckKicker{研究背景}
\SlideMeta{证据等级}{E3 / 草案/未批准规范}{来源状态：本地 PDF 已核验}
{\large\bfseries\color{Ink} TVM/Realm 使用设备时，需要把 device interface identity、DMA/MMIO、interrupt delivery 和 PCIe/CXL link 放到同一边界\par}\vspace{1.2mm}
\begin{columns}[T,totalwidth=\textwidth]
  \begin{column}{0.45\textwidth}
    \fcolorbox{Rule}{Paper}{%
  \begin{minipage}[t]{0.97\linewidth}
    \vspace{1.1mm}
    \hspace{1.4mm}\begin{minipage}{0.92\linewidth}
      \PanelTitle{内容说明}
      \scriptsize \begin{itemize}
  \item TVM/Realm 使用设备时，需要把 device interface identity、DMA/MMIO、interrupt delivery 和 PCIe/CXL link 放。
\end{itemize}
    \end{minipage}
    \vspace{1mm}
  \end{minipage}}
  \end{column}
  \begin{column}{0.49\textwidth}
    \fcolorbox{Rule}{Panel}{%
  \begin{minipage}[t]{0.97\linewidth}
    \vspace{1.1mm}
    \hspace{1.4mm}\begin{minipage}{0.92\linewidth}
      \PanelTitle{问题边界}
      \scriptsize {\tiny
\begin{tabularx}{\linewidth}{@{}YY@{}}
\toprule
\textbf{维度} & \textbf{说明} \\
\midrule
保护对象 & Confidential I/O protocol、trusted device interface、network endpoint \\
传统瓶颈 & 把 SPDM/TDISP/PCIe IDE/attested TLS 与 trusted device lifecycle 组合起来写。 \\
本文入口 & CoVE-IO and TDISP trusted-device lifecycle \\
证据限制 & E3 / 草案/未批准规范 \\
\bottomrule
\end{tabularx}
}
    \end{minipage}
    \vspace{1mm}
  \end{minipage}}
  \end{column}
\end{columns}
\SourceFootnote{来源：CoVE-IO and TDISP trusted-device lifecycle｜证据等级：E3 / 草案/未批准规范}
\end{frame}


\begin{frame}[t,shrink=6]{CoVE-IO and TDISP trusted-device lifecycle：解决方案}
\DeckKicker{解决方案}
\SlideMeta{证据等级}{E3 / 草案/未批准规范}{来源状态：本地 PDF 已核验}
{\large\bfseries\color{Ink} 用 CoVE-IO draft 描述 RISC-V confidential I/O lifecycle\par}\vspace{1.2mm}
\begin{columns}[T,totalwidth=\textwidth]
  \begin{column}{0.45\textwidth}
    \fcolorbox{Rule}{Paper}{%
  \begin{minipage}[t]{0.97\linewidth}
    \vspace{1.1mm}
    \hspace{1.4mm}\begin{minipage}{0.92\linewidth}
      \PanelTitle{内容说明}
      \scriptsize \begin{itemize}
  \item 用 CoVE-IO draft 描述 RISC-V confidential I/O lifecycle
  \item TDISP/TDISP XT 只作为公开 metadata 辅助，不推断 gated normative rules
\end{itemize}
    \end{minipage}
    \vspace{1mm}
  \end{minipage}}
  \end{column}
  \begin{column}{0.49\textwidth}
    \fcolorbox{Rule}{Panel}{%
  \begin{minipage}[t]{0.97\linewidth}
    \vspace{1.1mm}
    \hspace{1.4mm}\begin{minipage}{0.92\linewidth}
      \PanelTitle{核心设计拆解}
      \scriptsize \begin{itemize}
  \item 01. TDI/TDM/DSM and TSM/RDSM lifecycle
  \item 02. SPDM plus TDISP plus PCIe IDE composition
  \item 03. IOMMU/IOPMP/AIA trusted DMA and interrupt substrate
\end{itemize}
    \end{minipage}
    \vspace{1mm}
  \end{minipage}}
  \end{column}
\end{columns}
\SourceFootnote{来源：CoVE-IO and TDISP trusted-device lifecycle｜证据等级：E3 / 草案/未批准规范}
\end{frame}


\begin{frame}[t,shrink=6]{CoVE-IO and TDISP trusted-device lifecycle：实验结果}
\DeckKicker{实验结果}
\SlideMeta{证据等级}{E3 / 草案/未批准规范}{来源状态：本地 PDF 已核验}
{\large\bfseries\color{Ink} 规范/Survey，无新实验；CoVE-IO and TDISP trusted-device lifecycle 只支撑术语、taxonomy、接口语义或证据边界。\par}\vspace{1.2mm}
\begin{columns}[T,totalwidth=\textwidth]
  \begin{column}{0.45\textwidth}
    \fcolorbox{Rule}{Paper}{%
  \begin{minipage}[t]{0.97\linewidth}
    \vspace{1.1mm}
    \hspace{1.4mm}\begin{minipage}{0.92\linewidth}
      \PanelTitle{内容说明}
      \scriptsize \begin{itemize}
  \item 本页不展示独立系统实验，也不把二手归纳写成一手结果。
  \item 可支撑：RISC-V CoVE-IO draft maps SPDM and TDISP-style ideas into the RISC-V TEE-I/O track.
  \item 证据等级：E3 / Draft-not-ratified；来源状态：本地 PDF 已核验。
  \item 涉及具体平台、协议状态或数值时，转到原始论文、官方规范或厂商材料核验。
\end{itemize}
    \end{minipage}
    \vspace{1mm}
  \end{minipage}}
  \end{column}
  \begin{column}{0.49\textwidth}
    \fcolorbox{Rule}{Panel}{%
  \begin{minipage}[t]{0.97\linewidth}
    \vspace{1.1mm}
    \hspace{1.4mm}\begin{minipage}{0.92\linewidth}
      \PanelTitle{实验/证据边界}
      \scriptsize {\tiny
\begin{tabularx}{\linewidth}{@{}YYY@{}}
\toprule
\textbf{对象} & \textbf{可支撑} & \textbf{边界} \\
\midrule
数据/来源 & CoVE-IO and TDISP trusted-device lifecycle & 本地 PDF 已核验 \\
Baseline/对照 & 以论文或规范原文为准 & 不跨材料外推 \\
指标/对象 & 只使用当前来源可证明的信息 & E3 / 草案/未批准规范 \\
使用边界 & 可进入本方向演进图 & 不能替代更强证据 \\
\bottomrule
\end{tabularx}
}
    \end{minipage}
    \vspace{1mm}
  \end{minipage}}
  \end{column}
\end{columns}
\SourceFootnote{来源：CoVE-IO and TDISP trusted-device lifecycle｜证据等级：E3 / 草案/未批准规范}
\end{frame}


\begin{frame}[t,shrink=6]{CoVE-IO and TDISP trusted-device lifecycle：文章评价}
\DeckKicker{文章评价}
\SlideMeta{证据等级}{E3 / 草案/未批准规范}{来源状态：本地 PDF 已核验}
{\large\bfseries\color{Ink} 评价：CoVE-IO and TDISP trusted-device lifecycle 适合做 taxonomy/规范入口；规范/Survey，无新实验，不能替代一手系统证据。\par}\vspace{1.2mm}
\begin{columns}[T,totalwidth=\textwidth]
  \begin{column}{0.45\textwidth}
    \fcolorbox{Rule}{Paper}{%
  \begin{minipage}[t]{0.97\linewidth}
    \vspace{1.1mm}
    \hspace{1.4mm}\begin{minipage}{0.92\linewidth}
      \PanelTitle{内容说明}
      \scriptsize \begin{itemize}
  \item 设计优势：给 RISC-V TVM trusted I/O 提供最完整的公开 draft 架构入口。
  \item 局限性：v0.3.0 未 ratified，TDISP/XT 规范不可公开复核，不能写成 production CoVE-IO stack。
  \item 商业化潜力：指向 vendor NIC/DPU/accelerator assignment；风险在 SPDM证书、TDISP state、IDE key、ATS/PRI 和。
  \item 规范/Survey，无新实验；具体平台结论转到一手材料核验。
\end{itemize}
    \end{minipage}
    \vspace{1mm}
  \end{minipage}}
  \end{column}
  \begin{column}{0.49\textwidth}
    \fcolorbox{Rule}{Panel}{%
  \begin{minipage}[t]{0.97\linewidth}
    \vspace{1.1mm}
    \hspace{1.4mm}\begin{minipage}{0.92\linewidth}
      \PanelTitle{文章评价}
      \scriptsize {\tiny
\begin{tabularx}{\linewidth}{@{}YY@{}}
\toprule
\textbf{维度} & \textbf{判断} \\
\midrule
设计优势 & 给 RISC-V TVM trusted I/O 提供最完整的公开 draft 架构入口。 \\
主要局限 & v0.3.0 未 ratified，TDISP/XT 规范不可公开复核，不能写成 production CoVE-IO stack。 \\
商业落地 & 指向 vendor NIC/DPU/accelerator assignment；风险在 SPDM证书、TDISP state、IDE key... \\
讲稿定位 & 代表性改进; 证据强度 E3 \\
\bottomrule
\end{tabularx}
}
    \end{minipage}
    \vspace{1mm}
  \end{minipage}}
  \end{column}
\end{columns}
\SourceFootnote{来源：CoVE-IO and TDISP trusted-device lifecycle｜证据等级：E3 / 草案/未批准规范}
\end{frame}


\begin{frame}[t,shrink=6]{Separate but Together: TLS channel binding with remote attestation：内容摘要}
\DeckKicker{内容摘要}
\SlideMeta{证据等级}{E1 / Peer-reviewed endpoint channel-binding evidence}{来源状态：本地 PDF 已核验}
{\large\bfseries\color{Ink} TLS+RA binds an application secure channel to an attested TEE endpoint, anchoring endpoint/channel-binding...\par}\vspace{1.2mm}
\begin{columns}[T,totalwidth=\textwidth]
  \begin{column}{0.45\textwidth}
    \fcolorbox{Rule}{Paper}{%
  \begin{minipage}[t]{0.97\linewidth}
    \vspace{1.1mm}
    \hspace{1.4mm}\begin{minipage}{0.92\linewidth}
      \PanelTitle{内容说明}
      \scriptsize \begin{itemize}
  \item Binds an application secure channel to an attested TEE endpoint
  \item Anchors endpoint/channel-binding evidence
  \item Peer-reviewed endpoint/channel-binding evidence for integrating remote attestation into TLS.
\end{itemize}
    \end{minipage}
    \vspace{1mm}
  \end{minipage}}
  \end{column}
  \begin{column}{0.49\textwidth}
    \fcolorbox{Rule}{Panel}{%
  \begin{minipage}[t]{0.97\linewidth}
    \vspace{1.1mm}
    \hspace{1.4mm}\begin{minipage}{0.92\linewidth}
      \PanelTitle{证据定位}
      \scriptsize {\tiny
\begin{tabularx}{\linewidth}{@{}YY@{}}
\toprule
\textbf{字段} & \textbf{内容} \\
\midrule
论文定位 & 当前 SOTA/补充边界 \\
材料类型 & 系统论文 \\
证据等级 & E1 / Peer-reviewed endpoint channel-binding evidence \\
来源状态 & 本地 PDF 已核验 \\
\bottomrule
\end{tabularx}
}
    \end{minipage}
    \vspace{1mm}
  \end{minipage}}
  \end{column}
\end{columns}
\SourceFootnote{来源：Separate but Together: TLS channel binding with remote...｜证据等级：E1 / Peer-reviewed endpoint channel-binding evidence}
\end{frame}


\begin{frame}[t,shrink=6]{Separate but Together: TLS channel binding with remote attestation：研究背景}
\DeckKicker{研究背景}
\SlideMeta{证据等级}{E1 / Peer-reviewed endpoint channel-binding evidence}{来源状态：本地 PDF 已核验}
{\large\bfseries\color{Ink} SPDM and TDISP can identify device/control-plane state, but they do not prove where an application TLS...\par}\vspace{1.2mm}
\begin{columns}[T,totalwidth=\textwidth]
  \begin{column}{0.45\textwidth}
    \fcolorbox{Rule}{Paper}{%
  \begin{minipage}[t]{0.97\linewidth}
    \vspace{1.1mm}
    \hspace{1.4mm}\begin{minipage}{0.92\linewidth}
      \PanelTitle{内容说明}
      \scriptsize \begin{itemize}
  \item SPDM and TDISP identify device/control-plane state
  \item Application TLS needs binding to attestation evidence and workload identity
\end{itemize}
    \end{minipage}
    \vspace{1mm}
  \end{minipage}}
  \end{column}
  \begin{column}{0.49\textwidth}
    \fcolorbox{Rule}{Panel}{%
  \begin{minipage}[t]{0.97\linewidth}
    \vspace{1.1mm}
    \hspace{1.4mm}\begin{minipage}{0.92\linewidth}
      \PanelTitle{问题边界}
      \scriptsize {\tiny
\begin{tabularx}{\linewidth}{@{}YY@{}}
\toprule
\textbf{维度} & \textbf{说明} \\
\midrule
保护对象 & Confidential I/O protocol、trusted device interface、network endpoint \\
传统瓶颈 & 把 SPDM/TDISP/PCIe IDE/attested TLS 与 trusted device lifecycle 组合起来写。 \\
本文入口 & Separate but Together: TLS channel binding with remote... \\
证据限制 & E1 / Peer-reviewed endpoint channel-binding evidence \\
\bottomrule
\end{tabularx}
}
    \end{minipage}
    \vspace{1mm}
  \end{minipage}}
  \end{column}
\end{columns}
\SourceFootnote{来源：Separate but Together: TLS channel binding with remote...｜证据等级：E1 / Peer-reviewed endpoint channel-binding evidence}
\end{frame}


\begin{frame}[t,shrink=6]{Separate but Together: TLS channel binding with remote attestation：解决方案}
\DeckKicker{解决方案}
\SlideMeta{证据等级}{E1 / Peer-reviewed endpoint channel-binding evidence}{来源状态：本地 PDF 已核验}
{\large\bfseries\color{Ink} Integrate remote-attestation evidence into the TLS handshake or certificate path so the channel and TEE state...\par}\vspace{1.2mm}
\begin{columns}[T,totalwidth=\textwidth]
  \begin{column}{0.45\textwidth}
    \fcolorbox{Rule}{Paper}{%
  \begin{minipage}[t]{0.97\linewidth}
    \vspace{1.1mm}
    \hspace{1.4mm}\begin{minipage}{0.92\linewidth}
      \PanelTitle{内容说明}
      \scriptsize \begin{itemize}
  \item Integrate remote-attestation evidence into TLS
  \item Bind channel endpoint, TEE measurement, and verifier policy
\end{itemize}
    \end{minipage}
    \vspace{1mm}
  \end{minipage}}
  \end{column}
  \begin{column}{0.49\textwidth}
    \fcolorbox{Rule}{Panel}{%
  \begin{minipage}[t]{0.97\linewidth}
    \vspace{1.1mm}
    \hspace{1.4mm}\begin{minipage}{0.92\linewidth}
      \PanelTitle{核心设计拆解}
      \scriptsize \begin{itemize}
  \item 01. TLS endpoint and attestation binding
  \item 02. certificate/evidence composition
  \item 03. relying-party verification policy
\end{itemize}
    \end{minipage}
    \vspace{1mm}
  \end{minipage}}
  \end{column}
\end{columns}
\SourceFootnote{来源：Separate but Together: TLS channel binding with remote...｜证据等级：E1 / Peer-reviewed endpoint channel-binding evidence}
\end{frame}


\begin{frame}[t,shrink=6]{Separate but Together: TLS channel binding with remote attestation：实验结果}
\DeckKicker{实验结果}
\SlideMeta{证据等级}{E1 / Peer-reviewed endpoint channel-binding evidence}{来源状态：本地 PDF 已核验}
{\large\bfseries\color{Ink} USENIX ATC 2025 peer-reviewed system paper provides a TLS+RA prototype and evaluation.\par}\vspace{1.2mm}
\begin{columns}[T,totalwidth=\textwidth]
  \begin{column}{0.45\textwidth}
    \fcolorbox{Rule}{Paper}{%
  \begin{minipage}[t]{0.97\linewidth}
    \vspace{1.1mm}
    \hspace{1.4mm}\begin{minipage}{0.92\linewidth}
      \PanelTitle{内容说明}
      \scriptsize \begin{itemize}
  \item USENIX ATC 2025 peer-reviewed prototype and evaluation
  \item Supports endpoint/channel binding, not SPDM/TDISP device identity or DMA protection
\end{itemize}
    \end{minipage}
    \vspace{1mm}
  \end{minipage}}
  \end{column}
  \begin{column}{0.49\textwidth}
    \fcolorbox{Rule}{Panel}{%
  \begin{minipage}[t]{0.97\linewidth}
    \vspace{1.1mm}
    \hspace{1.4mm}\begin{minipage}{0.92\linewidth}
      \PanelTitle{实验/证据边界}
      \scriptsize {\tiny
\begin{tabularx}{\linewidth}{@{}YYY@{}}
\toprule
\textbf{对象} & \textbf{可支撑} & \textbf{边界} \\
\midrule
数据/来源 & Separate but Together: TLS channel binding with remote... & 本地 PDF 已核验 \\
Baseline/对照 & 以论文或规范原文为准 & 不跨材料外推 \\
指标/对象 & 只使用当前来源可证明的信息 & E1 / Peer-reviewed endpoint channel-binding evidence \\
使用边界 & 可进入本方向演进图 & 不能替代更强证据 \\
\bottomrule
\end{tabularx}
}
    \end{minipage}
    \vspace{1mm}
  \end{minipage}}
  \end{column}
\end{columns}
\SourceFootnote{来源：Separate but Together: TLS channel binding with remote...｜证据等级：E1 / Peer-reviewed endpoint channel-binding evidence}
\end{frame}


\begin{frame}[t,shrink=6]{Separate but Together: TLS channel binding with remote attestation：文章评价}
\DeckKicker{文章评价}
\SlideMeta{证据等级}{E1 / Peer-reviewed endpoint channel-binding evidence}{来源状态：本地 PDF 已核验}
{\large\bfseries\color{Ink} Fits the protocol/device-endpoint 主讲定位 better than FOLIO, which is a trusted-I/O contrast paper.\par}\vspace{1.2mm}
\begin{columns}[T,totalwidth=\textwidth]
  \begin{column}{0.45\textwidth}
    \fcolorbox{Rule}{Paper}{%
  \begin{minipage}[t]{0.97\linewidth}
    \vspace{1.1mm}
    \hspace{1.4mm}\begin{minipage}{0.92\linewidth}
      \PanelTitle{内容说明}
      \scriptsize \begin{itemize}
  \item 设计优势：Fits the protocol/device-endpoint 主讲定位 better than FOLIO, which is a trusted-I/O...
  \item 局限性：TLS+RA does not solve device assignment, TDISP state, IDE link protection, or...
  \item 商业化潜力：Useful for confidential service onboarding and key release; risks remain in...
  \item 证据边界：E1 / Peer-reviewed endpoint channel-binding evidence；本地 PDF 已核验。
\end{itemize}
    \end{minipage}
    \vspace{1mm}
  \end{minipage}}
  \end{column}
  \begin{column}{0.49\textwidth}
    \fcolorbox{Rule}{Panel}{%
  \begin{minipage}[t]{0.97\linewidth}
    \vspace{1.1mm}
    \hspace{1.4mm}\begin{minipage}{0.92\linewidth}
      \PanelTitle{文章评价}
      \scriptsize {\tiny
\begin{tabularx}{\linewidth}{@{}YY@{}}
\toprule
\textbf{维度} & \textbf{判断} \\
\midrule
设计优势 & Fits the protocol/device-endpoint 主讲定位 better than FOLIO, which is a... \\
主要局限 & TLS+RA does not solve device assignment, TDISP state, IDE link protection, or... \\
商业落地 & Useful for confidential service onboarding and key release; risks remain in... \\
讲稿定位 & 当前 SOTA/补充边界; 证据强度 E1 \\
\bottomrule
\end{tabularx}
}
    \end{minipage}
    \vspace{1mm}
  \end{minipage}}
  \end{column}
\end{columns}
\SourceFootnote{来源：Separate but Together: TLS channel binding with remote...｜证据等级：E1 / Peer-reviewed endpoint channel-binding evidence}
\end{frame}


\begin{frame}[t,shrink=6]{Confidential I/O protocol、trusted device interface、network endpoint：方向总结}
\DeckKicker{方向总结}
\SlideMeta{证据等级}{方向综述}{来源状态：公开材料综合}
{\large\bfseries\color{Ink} Confidential I/O protocol、trusted device interface、network endpoint 的结论是：三篇主讲材料共同给出技术演进，但仍要用证据等级约束每个结论。\par}\vspace{1.2mm}
\begin{columns}[T,totalwidth=\textwidth]
  \begin{column}{0.45\textwidth}
    \fcolorbox{Rule}{Paper}{%
  \begin{minipage}[t]{0.97\linewidth}
    \vspace{1.1mm}
    \hspace{1.4mm}\begin{minipage}{0.92\linewidth}
      \PanelTitle{内容说明}
      \scriptsize \begin{itemize}
  \item SPDM device identity and measurement：开创/基础入口，E0 / Spec-standard SOTA。
  \item CoVE-IO and TDISP trusted-device lifecycle：代表性改进，E3 / 草案/未批准规范。
  \item Separate but Together: TLS channel binding with remote...：当前 SOTA/补充边界，E1 / Peer-reviewed endpoint channel-binding evidence。
  \item 本方向结论：先用三篇主讲材料建立演进关系，再用证据等级控制机制、实验和商业化结论。
\end{itemize}
    \end{minipage}
    \vspace{1mm}
  \end{minipage}}
  \end{column}
  \begin{column}{0.49\textwidth}
    \fcolorbox{Rule}{Panel}{%
  \begin{minipage}[t]{0.97\linewidth}
    \vspace{1.1mm}
    \hspace{1.4mm}\begin{minipage}{0.92\linewidth}
      \PanelTitle{本方向方法对比}
      \scriptsize {\tiny
\begin{tabularx}{\linewidth}{@{}YYY@{}}
\toprule
\textbf{论文} & \textbf{定位} & \textbf{选择理由} \\
\midrule
SPDM device identity and measurement & 开创/基础入口 & SPDM is the first-party specification foundation for device identity... \\
CoVE-IO and TDISP trusted-device lifecycle & 代表性改进 & RISC-V CoVE-IO draft maps SPDM and TDISP-style ideas into the RISC-V... \\
Separate but Together: TLS channel binding with remote... & 当前 SOTA/补充边界 & Peer-reviewed endpoint/channel-binding evidence for integrating remote... \\
\bottomrule
\end{tabularx}
}
    \end{minipage}
    \vspace{1mm}
  \end{minipage}}
  \end{column}
\end{columns}
\SourceFootnote{来源：论文原文、官方规范或公开材料｜证据等级：见本页 badge}
\end{frame}
